\documentclass[11pt,letterpaper]{article}

\usepackage[margin=1in]{geometry}
\usepackage[T1]{fontenc}
\usepackage{lmodern}
\usepackage{microtype}
\usepackage{parskip}
\usepackage{enumitem}
\usepackage{booktabs}
\usepackage{longtable}
\usepackage{array}
\usepackage{titlesec}
\usepackage[hidelinks,colorlinks=true,linkcolor=black,urlcolor=black,citecolor=black]{hyperref}
\usepackage{xcolor}
\usepackage{fancyhdr}

\titleformat{\section}{\Large\bfseries}{\thesection}{1em}{}
\titleformat{\subsection}{\large\bfseries}{\thesubsection}{1em}{}
\titleformat{\subsubsection}{\normalsize\bfseries}{\thesubsubsection}{1em}{}

\pagestyle{fancy}
\fancyhf{}
\fancyhead[L]{\small The Theseus Guides --- IV. Architecture}
\fancyhead[R]{\small\thepage}
\renewcommand{\headrulewidth}{0.4pt}

\setlength{\parindent}{0pt}

\newcommand{\file}[1]{\texttt{#1}}
\newcommand{\mod}[1]{\texttt{#1}}

\title{\textbf{The Theseus Guides}\\[6pt]\large IV. Architecture:\\A Tour of the Repository and its Softwares}
\date{April 2026}
\author{}

\begin{document}
\maketitle
\thispagestyle{empty}

\begin{quote}\itshape
Theseus is a monorepo holding three principal pieces of software, two auxiliary web services, a research-experiment directory, a published-document catalogue, and the operational machinery to build, sign, and deploy them. This guide maps the territory.
\end{quote}

\vspace{1em}
\tableofcontents
\newpage

% =======================================================================
\section{The repository at a glance}
% =======================================================================

The repository is organized so that each software artifact lives in its own directory, with shared \file{docs/}, \file{scripts/}, \file{packages/}, and \file{deploy/} folders for cross-cutting concerns. The top-level \file{README.md} is the entry point for new contributors; it lists the downloadable installers, the live Codex URL, and the guiding principle.

\begin{longtable}{p{3.4cm} p{11.1cm}}
\toprule
\textbf{Directory} & \textbf{Purpose} \\
\midrule
\endhead
\file{noosphere/}            & The core epistemic engine (Python). Ingests transcripts and essays, extracts claims, runs six-layer coherence, builds the knowledge graph, synthesizes conclusions, runs adversarial review and calibration, exposes a CLI. Packaged as a standalone binary via PyInstaller. \\
\file{dialectic/}            & Live-conversation analyzer (PyQt6 desktop app). Streams microphone audio through Whisper, segments into claims, flags contradictions on a dashboard in real time. Distributed as macOS DMG and Windows NSIS installer. \\
\file{theseus-codex/}        & Founders' web control plane and Electron desktop app (Next.js 16 / React 19 / Prisma 7). Authentication, artifact upload, ingestion dispatch, conclusion review, adversarial page, decay dashboard, rigor-gate queue. \\
\file{theseus-public/}       & Publication-only public website (Next.js static export). Signed method docs, MIP registry, per-conclusion provenance, methodology refusal dashboard. \\
\file{researcher\_api/}      & Rate-limited HTTP API for outside investigators --- claim extraction, coherence, embedding, calibration, replay on researcher-supplied text only. \\
\file{ideologicalOntology/}  & Research experiments in contradiction geometry, including the Embedding Geometry Conjecture validation suite and the Reverse Marxism framework. Data and analysis scripts, not deployed software. \\
\file{dialectic/}            & (see above) \\
\file{docs/}                 & Published PDFs (research, architecture, product) plus operational markdowns and per-area specs. \\
\file{reference/}            & Starting-material snapshots (git-ignored; preserved for provenance). \\
\file{packages/}             & Shared TypeScript packages (\texttt{theseus-api-types}) used by the web frontends. \\
\file{deploy/}               & Docker Compose stacks and Kubernetes Helm charts for the web services. \\
\file{scripts/}              & Shared build infrastructure: macOS code signing, Apple notarization, Windows Authenticode, invariant-check lint scripts. \\
\file{.github/}              & CI/CD workflows: cross-platform desktop builds, Docker image builds, release orchestration, round-3 invariant checks, red-team workflows. \\
\file{Podcast talks/}        & Raw source transcripts of the founders' weekly four-hour discussions. \\
\file{noosphere\_data/}      & On-disk state for a local Noosphere instance (embeddings, graph, registries). Git-ignored. \\
\bottomrule
\end{longtable}

Two details are architecturally important. First, the separation between \file{theseus-codex/} (founders only) and \file{theseus-public/} (publication only, read-only) is not a convention --- it is enforced at the storage layer: writes to the public content store that bypass the rigor gate are rejected by an assertion, and a CI lint verifies that every publication handler wears the \texttt{@gated} decorator. Second, every method invocation anywhere in the system writes a signed ledger entry; this is not optional instrumentation, it is a pre-hook on the method registry.

% =======================================================================
\section{Noosphere --- the epistemic engine}
% =======================================================================

Noosphere is the Brain of the Firm. It is a Python package at \file{noosphere/noosphere/} containing approximately 40 top-level modules and 15 subpackages, with \textasciitilde 89 Python modules and 95+ test files overall. It is not a monolith; it is a composition of subsystems with clean cross-boundary interfaces.

\subsection{What it does, from ten thousand feet}

Noosphere takes the weekly four-hour podcast transcripts (plus founder essays, memos, and Dialectic session exports) as input, decomposes them into atomic claims, evaluates those claims across multiple coherence methods, tracks principles and their evolution over time, and synthesizes firm-tier conclusions with calibrated confidence. Every judgment it makes is a registered \emph{method} routed through a central registry; every method invocation is recorded in the append-only ledger; every evidence flow between claims and conclusions is represented as a typed cascade graph.

\subsection{The flat modules --- the core pipeline}

The top of \file{noosphere/noosphere/} holds single-file modules implementing the foundational pipeline. Many have been ported into the method registry as registered methods; the flat-module originals remain as the canonical implementations.

\begin{longtable}{p{3.6cm} p{10.9cm}}
\toprule
\textbf{Module} & \textbf{Responsibility} \\
\midrule
\endhead
\mod{models.py}             & Pydantic types for every cross-boundary object (\texttt{Artifact}, \texttt{Chunk}, \texttt{Claim}, \texttt{Principle}, \texttt{Conclusion}, and all round-3 types: \texttt{Method}, \texttt{MethodInvocation}, \texttt{LedgerEntry}, \texttt{CascadeNode}/\texttt{Edge}, \texttt{DecayPolicy}, \texttt{ReviewReport}, \texttt{RigorVerdict}). \\
\mod{store.py}              & SQLModel-backed persistence; typed CRUD; chain-integrity enforcement for ledger entries and cascade edges. \\
\mod{ids.py}                & Deterministic content-addressed identifiers for every domain object. \\
\mod{config.py}             & Pydantic-settings configuration, reading environment variables and optional \file{theseus.toml}. \\
\mod{observability.py}      & Structlog JSON logger with per-run correlation IDs. \\
\mod{llm.py}                & LLM client protocol plus Anthropic and OpenAI implementations plus a mock for tests. \\
\mod{embeddings.py}         & Embedding client protocol plus sentence-transformers implementation plus on-disk cache. \\
\mod{ingester.py}           & Markdown, plain-text, and WebVTT transcript ingest with YAML frontmatter, speaker-turn segmentation, and Dialectic session import. \\
\mod{claim\_extractor.py}   & Chunk-to-claim extraction via LLM with structured JSON output; cached per chunk. \\
\mod{classifier.py}         & Zero-shot MNLI verification of claim type; flags disagreements between LLM-assigned and classifier-assigned labels. \\
\mod{embed\_pass.py}        & Orchestrated embedding pass over claims lacking embeddings; batched; supports rebuild. \\
\mod{ontology.py}           & Growing topic ontology via HDBSCAN on UMAP-reduced embeddings; persistent cluster identities; spaCy-based entity linking. \\
\mod{geometry.py}           & Contradiction-geometry detection via Householder reflection and difference-vector sparsity per the Embedding Geometry Conjecture. \\
\mod{conclusions.py}        & Synthesizes firm-tier, founder-tier, and open-tier conclusions with tiered confidence discount. \\
\mod{synthesis.py}          & End-to-end synthesis pipeline tying ingest, extraction, coherence, and conclusion into a single runner. \\
\mod{adversarial.py}        & Generates structured objections to firm-level conclusions, runs them through the six-layer coherence stack, and persists verdicts. \\
\mod{redteam.py}            & Red-team prompts and harness for adversarial robustness testing. \\
\mod{meta\_analysis.py}     & Five-criterion meta-analysis gate over candidate conclusions (severity, coherence, freshness, provenance, calibration). \\
\mod{predictive\_extractor.py} & Extracts falsifiable predictions with resolution criteria. \\
\mod{resolution.py}         & Resolves predictions against observed outcomes. \\
\mod{scoring.py}            & Calibration-scoreboard computation (Brier, log loss, ECE). \\
\mod{retrieval.py}          & Embedding-based corpus retrieval for downstream methods and the research advisor. \\
\mod{founders.py}, \mod{voices.py}           & Voice decomposition --- separates named founder positions from firm-tier claims. \\
\mod{literature.py}         & External-literature search and claim-match scoring against the firm's corpus. \\
\mod{research\_advisor.py}  & Generates ranked research suggestions based on open questions and blind spots. \\
\mod{temporal.py}, \mod{temporal\_replay.py} & Principle emergence, drift, and strengthening; counterfactual temporal replay. \\
\mod{tenancy.py}            & Multi-tenant partitioning for cloud deployments. \\
\mod{storage\_client.py}    & Object-store abstraction for ledger-referenced input/output blobs. \\
\mod{backup\_restore.py}    & Archive-and-restore of SQLite plus data directory. \\
\mod{orchestrator.py}       & Coordinates multi-step pipelines. \\
\mod{frozen\_support.py}    & Path resolution for PyInstaller frozen mode. \\
\mod{updater.py}            & Background version checker against a release manifest URL. \\
\mod{typer\_cli.py}         & Primary Typer CLI entry point with subcommand auto-discovery. \\
\bottomrule
\end{longtable}

\subsection{The coherence subpackage}

\file{noosphere/noosphere/coherence/} holds the six-layer coherence stack --- Theseus's principal contribution to how one measures belief-set consistency. Each layer is independent, has clear empirical semantics, and produces output the aggregator can combine.

\begin{longtable}{p{3.1cm} p{11.4cm}}
\toprule
\textbf{Layer} & \textbf{What it measures} \\
\midrule
\endhead
\mod{nli.py}               & Natural-language inference verdicts (entail / neutral / contradict probabilities) via a DeBERTa-v3 MNLI/ANLI head, batched. \\
\mod{argumentation.py}     & Formal argumentation-framework analysis (Dung-style admissibility, preferred / grounded extensions). \\
\mod{probabilistic.py}     & Probabilistic-consistency check: is the claim set jointly satisfiable under any coherent probability assignment? \\
\mod{geometry.py}          & Embedding-geometry test using the validated Embedding Geometry Conjecture. \\
\mod{information.py}       & Information-theoretic compressibility: claim-set entropy and redundancy. \\
\mod{judge.py}             & LLM-as-judge with structured output and disagreement-with-other-layers flagging. \\
\mod{aggregator.py}        & Combines the six layers into a \texttt{SixLayerScore} with explicit \texttt{unresolved} output when layers disagree materially. \\
\mod{engine.py}            & Orchestrates a coherence run end-to-end; populates the coherence cache. \\
\mod{scheduler.py}         & Picks pairs to score, prioritizing load-bearing-for-conclusions pairs. \\
\mod{calibration.py}       & Maintains calibration bundles of reference pairs with gold labels for model-upgrade recalibration. \\
\mod{metrics.py}           & Macro-F1 and per-layer agreement metrics for regression testing. \\
\mod{cache.py}             & On-disk cache of per-pair verdicts, invalidated on model-version change. \\
\bottomrule
\end{longtable}

The aggregator is deliberately conservative: when layers disagree past a threshold, the output is \texttt{unresolved} rather than a point verdict. \texttt{unresolved} is a first-class output in every downstream consumer, not an error state.

\subsection{Subpackages}

Beyond the flat core and the coherence stack, Noosphere has fifteen subpackages, each owning a slice of the system:

\begin{longtable}{p{2.4cm} p{12.1cm}}
\toprule
\textbf{Package} & \textbf{Role} \\
\midrule
\endhead
\file{methods/}       & Central registry with the \texttt{@register\_method} decorator, pre/post-hook API, legacy parity-test copies, and ported implementations (nli\_scorer, contradiction\_geometry, six\_layer\_coherence, suggest\_research, extract\_claims, classify\_claim\_type, synthesize\_conclusion, decompose\_voice, external\_claim\_match, extract\_from\_corpus, extract\_prediction). \\
\file{ledger/}        & Append-only, Ed25519-signed, Merkle-chained audit log. Every method invocation produces one entry. Externally verifiable and exportable as an audit bundle. \\
\file{cascade/}       & Typed evidence-flow graph linking artifacts $\to$ chunks $\to$ claims $\to$ principles $\to$ conclusions with eleven relation kinds. Supports proof-bundle export, cut-reports, single-point-of-failure diagnostics, cycle detection, and rendering. \\
\file{evaluation/}    & Counterfactual calibration rig. \texttt{CorpusSlicer} presents a read-only view pinned to a past \texttt{as\_of}; methods invoked against it raise \texttt{EmbargoViolation} on out-of-slice reads. Computes Brier, log loss, ECE, reliability curves per method. \\
\file{inference/}     & Inverse inference engine. Given a resolved event, back-propagates to corpus claims that would have predicted it; produces a structured blind-spot report. \\
\file{external\_battery/} & Monthly benchmarks against external corpora (Good Judgment, Metaculus, ClaimReview, replication projects) with a five-way failure taxonomy. \\
\file{peer\_review/}  & Adversarial peer-review swarm: seven declared reviewer roles (methodological, evidential, statistical, adversarial-literature, replication, rhetorical, humility), reviewer-level calibration discounting. \\
\file{decay/}         & Policy-driven re-validation. Policies: fixed-interval, evidence-changed, method-version-bump, embedding-drift, outcome-observed, calibration-regression. Freshness is derived, not cached. \\
\file{rigor\_gate/}   & Publication choke point. Nine declared checks; failures block; founder overrides are public and ledger-recorded. Refusal rates published monthly. \\
\file{transfer/}      & Packages a registered method as a Docker-reproducible, signed artifact with spec, rationale, implementation, adapter, eval card, and transfer-degradation studies. \\
\file{docgen/}        & Compiles signed, versioned \texttt{MethodDoc} bundles; AST-diff changelog enforcement on version bumps. \\
\file{interop/}       & Builds Methodology Interoperability Packages (MIPs). Ships with a minimal non-Turing-complete YAML workflow language. \\
\file{mitigations/}   & Security and robustness modules: ingestion guard, citation anchor, tenant audit, calibration guard, temporal guard, embedding-text validation, coherence judge guard. \\
\file{cli\_commands/} & Typer subcommand modules auto-discovered by the main CLI. \\
\bottomrule
\end{longtable}

\subsection{Distribution}

Noosphere is packaged as a standalone CLI binary via PyInstaller. The \file{noosphere.spec} bundles Alembic migration scripts and all Python dependencies into a single-directory distribution. ML model weights (sentence-transformers, spaCy) are downloaded on first run rather than bundled. Platform build scripts produce a macOS \file{.tar.gz} and a Windows NSIS installer that adds Noosphere to the system PATH.

% =======================================================================
\section{Dialectic --- the live conversation analyzer}
% =======================================================================

Dialectic listens to a discussion while it is happening. The directory \file{dialectic/} holds a small PyQt6 application (\textasciitilde 11 modules, 5 test files) designed to surface the logical structure of a conversation in real time to the people in the room.

\subsection{The pipeline}

\file{audio.py} opens a microphone stream via \file{sounddevice}, runs voice-activity detection to cut audio into speech segments, and pushes segments into a thread-safe queue. \file{transcriber.py} consumes the queue with \file{faster-whisper} (CTranslate2-backed) and emits partial and final transcript tokens. \file{analyzer.py} takes finalized utterances, extracts claims with speaker attribution, embeds them with sentence-transformers MiniLM, tracks topic clusters, and runs a DeBERTa-v3 NLI head pairwise across recent claims to detect contradictions above a configurable threshold. \file{dashboard.py} renders all of this as a Qt interface with three panes: running transcript, live claim graph, and contradiction alerts. The design goal is a latency under two seconds from speech to claim appearance on a laptop GPU or Apple Silicon.

\subsection{The live interlocutor (SP09)}

A distinguishing feature: Dialectic can surface short, third-person prompts into the conversation itself when participants opt in per session. Modes range from silent (default) through passive overlay, conversational with optional TTS after a pause, and tutor (higher rate, requires explicit acknowledgment). Each candidate prompt passes a conservative quality gate and a per-session budget. Every intervention and every drop is logged to JSONL with a reflection bundle for post-session review in the Theseus Codex.

\subsection{Pairing with Noosphere}

Saved sessions are JSON Lines, one line per finalized claim with timestamp, speaker, text, embedding, and any contradictions found. These files are the native input format for Noosphere's ingester. When environment variables \texttt{DIALECTIC\_CLOUD\_URL} and \texttt{DIALECTIC\_CLOUD\_API\_KEY} are set, finalized transcripts and reflection bundles auto-upload to the Codex's \file{/api/upload} endpoint after each session ends; audio \file{.wav} files stay local as provenance.

\subsection{Distribution}

Dialectic is packaged as a standalone desktop application via PyInstaller. The \file{dialectic.spec} bundles PyQt6, sounddevice, faster-whisper, sentence-transformers, and torch. Build scripts produce a macOS \file{.app} bundle in a DMG with a drag-to-Applications layout, and a Windows NSIS installer with Start Menu and Desktop shortcuts. A background update checker (\file{updater.py}) polls a version manifest and notifies the user within the Qt dashboard when a newer release is available.

% =======================================================================
\section{Theseus Codex --- the founders' control plane}
% =======================================================================

The Theseus Codex at \file{theseus-codex/} is where founders operate Noosphere without touching the command line. The stack is Next.js 16 (App Router) on React 19 with Prisma 7 over SQLite, plus an Electron 35 desktop wrapper. Live at \url{https://mqtheseuswork-qiw6.vercel.app}.

\subsection{The seven-tab dashboard}

The top navigation has seven destinations. Four are group labels that expand sub-navigation when clicked:

\begin{longtable}{@{}p{2.6cm} p{1.7cm} p{10.0cm}@{}}
\toprule
\textbf{Top nav} & \textbf{Default} & \textbf{What lives inside} \\
\midrule
\endhead
Dashboard   & \file{/dashboard}      & Single page --- recent uploads, new conclusions, drift events. \\
Upload      & \file{/upload}         & Single page --- drag-and-drop ingestion form. \\
Conclusions & \file{/conclusions}    & List of firm claims; click for a four-tab detail view (Overview, Provenance, Cascade, Peer review). \\
Review      & \file{/contradictions} & Five sub-tabs: Contradictions, Open questions, Adversarial, Layer review, Calibration. \\
Library     & \file{/voices}         & Four sub-tabs: Voices, Literature, Reading queue, Research. \\
Publication & \file{/publication}    & Internal review queue plus static-site export. \\
Ops         & \file{/provenance}     & Seven sub-tabs: Provenance, Eval runs, Post-mortem, Decay, Rigor gate, Methods, Founders. \\
\bottomrule
\end{longtable}

Every page carries a consistent banner with its name, a one-sentence description of what it shows, and a one-sentence description of when a founder would use it. This banner is the orientation layer; nothing below it makes sense if skipped.

\subsection{How the Codex talks to Noosphere}

The Codex does not reason; it renders. Almost every page is a view over a specific table in the underlying Postgres (production) or SQLite (development / Electron) database. Operationally, the portal points at the same database file Noosphere writes, so the \texttt{adversarial\_challenge}, \texttt{conclusion}, and related tables are shared directly; there is no API marshalling layer for internal reads. Mutations that write to publicly-visible state go through a Python bridge (\file{src/lib/ingest.ts}) that shells out to the Noosphere CLI under the portal user's environment, streaming stdout back into the database for live progress.

\subsection{Distribution}

The Codex ships as both a web service and a desktop application. The Electron wrapper in \file{theseus-codex/desktop/} (six TypeScript modules: \file{main.ts}, \file{next-server.ts}, \file{db-path.ts}, \file{preload.ts}, \file{updater.ts}, \file{dev.ts}) embeds the Next.js standalone build, resolves the SQLite path to platform-appropriate user-data directories, runs Prisma migrations on launch, and uses \texttt{electron-updater} to check GitHub Releases. The \file{electron-builder.yml} produces a macOS DMG (universal architecture, hardened runtime) and a Windows NSIS installer. For the web path, a multi-stage Dockerfile builds the Next.js standalone output, generates the Prisma client, pre-runs migrations to produce a seed database, and serves from \texttt{node:20-alpine}.

% =======================================================================
\section{Theseus Public --- the publication-only site}
% =======================================================================

The public site at \file{theseus-public/} is the firm's outward intellectual face. It reads only from the publication store; it has no write paths. Each conclusion published here carries a provenance block linking to the ledger entries that produced it, a sanitized adversarial history, and the corpus hash at the time of publication.

The site hosts:

\begin{itemize}[leftmargin=1.5em,itemsep=0.1em]
\item \file{/methods/} --- the public registry of packaged, signed method documentation bundles.
\item \file{/interop/} --- published Methodology Interoperability Packages (MIPs) for external adoption.
\item \file{/methodology/rigor/} --- monthly pass/fail rates of the publication rigor gate.
\item \file{/methodology/overrides/} --- every founder override of the rigor gate, with full justification.
\item \file{/methodology/decay/} --- aggregated statistics on what was retired and why.
\item \file{/open-questions/} and \file{/responses/} --- active unresolved topics and structured replies to external critiques.
\end{itemize}

Writes that bypass the rigor gate are architecturally impossible: a storage-layer assertion rejects un-gated writes, and a CI lint verifies that every publication handler wears the \texttt{@gated} decorator. Distribution is via Next.js static export, served through \texttt{nginx:1.27-alpine} in a multi-stage Docker image, with Vercel and Netlify configurations as alternatives.

% =======================================================================
\section{Researcher API --- the outside-investigator surface}
% =======================================================================

The Researcher API at \file{researcher\_api/} is a rate-limited, authenticated HTTP service for academic or industrial researchers who want to apply Theseus methods to their own data without access to the firm's internal corpus. It exposes claim extraction, coherence scoring, embedding, calibration scoring, and a stub replay endpoint --- all restricted to researcher-supplied text.

It is a separate FastAPI service with its own API-key registry, privacy policy, security documentation, acceptable-use policy, and citation guidance. The split from Noosphere is architectural: the Researcher API cannot read firm-internal data and does not touch the shared store.

% =======================================================================
\section{Ideological Ontology --- research experiments}
% =======================================================================

\file{ideologicalOntology/} holds the empirical work that validates the firm's methodological claims. It contains the Embedding Geometry Conjecture validation suite (covered in Guide III), the Reverse Marxism ideology-flipping framework, a catalogue of nine contradiction-geometry experiments with methods and results, and four written documents (\file{Coherence\_Formula\_Walkthrough.pdf}, \file{Contradiction\_Detection\_Guide.pdf}, \file{Domain\_Relative\_Coherence.md}, \file{Formal\_Mathematical\_Framework.pdf}).

This is not deployed software. It is the experimental substrate that justifies why the coherence engine is shaped the way it is. When an implementation decision in \file{noosphere/} needs justification, the trail typically runs from a production module, through a written document in \file{docs/}, to an experiment in \file{ideologicalOntology/}, which is where the empirical backing lives.

% =======================================================================
\section{Operational and deployment infrastructure}
% =======================================================================

\subsection{CI/CD}

\file{.github/workflows/} contains six GitHub Actions workflows: \file{build-dialectic.yml}, \file{build-noosphere.yml}, \file{build-theseus-codex.yml}, \file{build-docker.yml}, \file{release.yml}, and \file{round3.yml} for architectural invariant checks plus \file{noosphere-redteam.yml} for adversarial/security testing. Desktop builds run on matrix macOS-14 (ARM) and Windows-latest. Dependabot runs weekly against npm, pip, and GitHub Actions. The release workflow is triggered on \texttt{v*} tags and produces a draft GitHub Release via \texttt{softprops/action-gh-release}.

\subsection{Code signing and notarization}

Three shared scripts in \file{scripts/} handle cross-platform signing. \file{codesign\_macos.sh} deep-signs all nested binaries with a Developer ID certificate, sets hardened runtime and timestamp, and gracefully skips when no signing identity is provided. \file{notarize\_macos.sh} submits DMG or ZIP to Apple's notarization service via \texttt{xcrun notarytool}, waits up to 600 seconds, and staples the ticket on success. \file{codesign\_windows.ps1} signs all \texttt{.exe} and \texttt{.dll} files with an Authenticode certificate via \texttt{signtool}, using SHA-256 and DigiCert timestamping.

All three scripts are designed for CI integration: they exit 0 with a warning when credentials are not configured, so unsigned development builds still succeed.

\subsection{Deployment}

\file{deploy/compose/} contains the production Docker Compose configuration. \file{deploy/helm/theseus/} contains a Kubernetes Helm chart with templates for portal and ingest-worker deployments. The root-level \file{docker-compose.yml} runs the Codex on port 3000 and the Public site on port 8080 with a shared \file{portal-data} volume for SQLite persistence. \file{theseus-public/vercel.json} and \file{theseus-public/netlify.toml} configure managed hosting with appropriate security headers.

% =======================================================================
\section{Who uses what}
% =======================================================================

Three roles interact with the system, each through a distinct surface.

\textbf{Founder.} The primary user. Works through the Theseus Codex (web UI or Electron desktop app, credentialed) and secondarily through the CLI when diagnosing. Uploads transcripts and essays, reviews synthesized conclusions, triages contradictions and adversarial challenges, authorizes rigor-gate overrides, retires invalidated objects. Every founder action is ledger-recorded under an \texttt{Actor(kind="human")}.

\textbf{Researcher (external).} An outside academic or industrial user. Accesses the Researcher API via HTTP with an issued API key. Can submit their own text to claim extraction, coherence scoring, embedding, calibration, and replay endpoints. Cannot touch firm-internal data. Can download packaged methods and MIPs from the public site and re-run them locally against their own corpus.

\textbf{Public reader.} Anyone on the internet. Reads the public site: published conclusions with full provenance, adversarial histories, method documentation, MIP registry, methodology refusal dashboards, open questions, and the firm's structured responses to external critiques. Cannot write.

\subsection{The operational surface for the firm itself}

The firm's own operational surface is the Typer CLI and the Codex. Representative CLI verbs:

\begin{itemize}[leftmargin=1.5em,itemsep=0.1em]
\item \texttt{noosphere ingest} / \texttt{synthesize} / \texttt{research} --- core pipeline commands.
\item \texttt{noosphere methods} --- list, show, run, diff, package, document, extract-candidates.
\item \texttt{noosphere ledger} --- verify the chain, export an audit bundle.
\item \texttt{noosphere cascade} --- explain, cut, export a proof bundle, diagnostics.
\item \texttt{noosphere eval counterfactual / external} --- run calibration and external-battery benchmarks.
\item \texttt{noosphere inverse} --- post-mortem on a resolved event.
\item \texttt{noosphere review} --- swarm, rebuttals-pending, reviewer-calibration.
\item \texttt{noosphere decay} --- status, revalidate, schedule, retire.
\item \texttt{noosphere interop} --- build, run, scaffold, submit-transfer.
\item \texttt{noosphere gate} --- submit, status, override, refusal-report.
\item \texttt{noosphere backup} / \texttt{restore} --- archive and restore full state.
\item \texttt{noosphere adversarial} --- generate and evaluate adversarial challenges.
\end{itemize}

% =======================================================================
\section{The technology stack, summarized}
% =======================================================================

\begin{longtable}{p{3cm} p{11.5cm}}
\toprule
\textbf{Layer} & \textbf{Technologies} \\
\midrule
\endhead
Backend/Engine & Python 3.11: sentence-transformers, transformers, spaCy, scikit-learn, HDBSCAN, UMAP, NetworkX, PyNaCl (Ed25519), zstandard, z3-solver. SQLModel + Alembic over SQLite or PostgreSQL. LLM clients: Anthropic Claude, OpenAI. \\
Frontends & Next.js 16 (App Router), React 19, TypeScript 5, TailwindCSS 4. Prisma 7 ORM with better-sqlite3. \\
Desktop & Electron 35 with electron-builder 26 and electron-updater (Theseus Codex). PyQt6 with pyqtgraph and qasync (Dialectic). PyInstaller for all Python apps. \\
Audio/NLP & faster-whisper (CTranslate2), sounddevice (PortAudio), DeBERTa-v3 NLI, sentence-transformers MiniLM. \\
Deployment & Docker multi-stage (node:20-alpine, nginx:1.27-alpine), Docker Compose, Kubernetes Helm charts. Vercel/Netlify for static hosting. \\
CI/CD & GitHub Actions on macOS 14 (ARM) and Windows-latest. Code signing: macOS \texttt{codesign} plus \texttt{notarytool}; Windows \texttt{signtool} Authenticode. Dependabot. \\
Testing & pytest (Python), Vitest (TypeScript), invariant-check scripts, red-team harness, calibration regression suites. \\
\bottomrule
\end{longtable}

The stack is deliberately conservative. No exotic databases, no bleeding-edge ML frameworks that change their API on a semi-monthly cadence, no cloud-specific primitives that lock the firm into a single vendor. Where the firm is novel is in \emph{what it does} with a boring stack, not in the stack itself.

\end{document}
